% R911: gerado; alterar dados e gerador, não esta prancha.

\refstepcounter{figure}\label{fig:AR09-mental}
\noindent{\large\sffamily\bfseries\color{navy}AR09 · Mapa mental}\par\smallskip
\noindent\textbf{Confiança, economia e governança}\par
\noindent\textit{Quem pode receber a tarefa e sob que regras?}\par\medskip
\begin{center}

\tikzset{
 atlasnode/.style={draw=steel,fill=blueBg,rounded corners=2mm,align=center,
  text width=4cm,inner sep=5pt,font=\fontsize{9.5}{11.8}\selectfont,text=textmain},
 atlasroot/.style={atlasnode,fill=navy,text=white,text width=2.8cm,font=\fontsize{10}{12}\selectfont\bfseries},
 atlasdecision/.style={diamond,aspect=2.1,draw=amber,fill=amberBg,align=center,
  text width=3cm,inner sep=2pt,font=\fontsize{9.5}{11.5}\selectfont},
 atlascall/.style={-Stealth,draw=steel,thick},
 atlasevent/.style={-Stealth,draw=violet,dashed,thick},
 atlasdata/.style={-Stealth,draw=teal,densely dotted,thick},
 atlasstructure/.style={draw=slate,thick},
 atlasboundary/.style={-Stealth,draw=amber,dash dot,thick},
 atlasedgelabel/.style={fill=papel,inner sep=1.3pt,font=\fontsize{9}{10}\selectfont}
}

\begin{tikzpicture}
\node[atlasroot] (root) at (0,0) {AR09\\Confiança, economia e governança};

\node[atlasnode] (branch0) at (-4.5,2.5) {\textbf{Elegibilidade}\\capacidade\\confiança\\limiar};
\draw[atlasstructure] (root) -- (branch0);

\node[atlasnode] (branch1) at (4.5,2.5) {\textbf{Decisão}\\risco\\negação\\histórico};
\draw[atlasstructure] (root) -- (branch1);

\node[atlasnode] (branch2) at (-4.5,-2.5) {\textbf{Economia interna}\\saldo\\reserva\\liquidação};
\draw[atlasstructure] (root) -- (branch2);

\node[atlasnode] (branch3) at (4.5,-2.5) {\textbf{Governança}\\rastro\\resultado\\supervisão};
\draw[atlasstructure] (root) -- (branch3);
\end{tikzpicture}
\end{center}
{\small\textbf{Figura \thefigure.} Mapa mental --- Confiança, economia e governança}\par\smallskip
\noindent\textbf{Como funciona e por quê.} Parta do centro e leia os quatro ramos. As linhas associam conceitos; não representam chamadas de funções.\par
\noindent\textbf{Justificativa.} No caminho de propostas, o primário consulta TrustEngine.execute e exclui decisões negadas ou com erro. O gate usa o histórico da ação. Após escolher um agente, o primário tenta reservar recursos na economia interna. O relato de conclusão atualiza o aprendizado e a liquidação conforme o resultado registrado.\par
\noindent\textbf{Limite.} O agrupamento é didático; cada responsabilidade será localizada no diagrama e conferida no fluxo.\par\smallskip
\fontecore{marceloclaro/orchestrator.py}{MarceloClaroOrchestrator._on_cfp}
\fontecore{marceloclaro/orchestrator.py}{MarceloClaroOrchestrator.report_completion}
